\documentclass[11pt]{article}
\usepackage{amsmath, amssymb}
\usepackage{geometry}
\usepackage{hyperref}
\usepackage{xcolor}
\geometry{margin=2.5cm}

\title{Finite-Size Majorana Splitting in the Kitaev Chain:\\
Why do the ``spikes'' in the finite-size spectrum worsen with larger $L$?}
\author{Seminar notes}
\date{\today}

\begin{document}
\maketitle

\section{Setup}

We study the Kitaev chain at $t = \Delta = 1$ with open boundary conditions (OBC).
The BdG Hamiltonian is the real-symmetric matrix
\begin{equation}
  M = \begin{pmatrix} h & d \\ -d & -h \end{pmatrix}, \qquad
  h_{ij} = -\mu\delta_{ij} - t(\delta_{i,j+1}+\delta_{i+1,j}), \qquad
  d_{ij} = \Delta(\delta_{i,j+1}-\delta_{i+1,j}).
\end{equation}
Particle-hole symmetry (PHS) guarantees eigenvalues come in exact $\pm E$ pairs.
Plot 4 tracks the \emph{six smallest positive eigenvalues} as $\mu$ varies.

\section{The ideal edge-mode energy}

\subsection{Zero-energy condition in the complex BZ}

To find the edge-mode energy, we look for a zero-energy solution that decays
into the bulk. Set $E = 0$ and try an Ansatz $\psi_j \sim \lambda^j$.
Substituting into the bulk BdG equations gives the characteristic equation
\begin{equation}
  \bigl(-\mu - t(\lambda + \lambda^{-1})\bigr)^2 + \Delta^2(\lambda - \lambda^{-1})^2 = 0.
\end{equation}
For $t = \Delta$ this factors to
\begin{equation}
  \bigl(-\mu - t(\lambda + \lambda^{-1}) + \Delta(\lambda - \lambda^{-1})\bigr)
  \bigl(-\mu - t(\lambda + \lambda^{-1}) - \Delta(\lambda - \lambda^{-1})\bigr) = 0.
\end{equation}
The two branches give the solutions
\begin{equation}
  \lambda_A = \frac{\mu}{2t}, \qquad \lambda_B = \frac{2t}{\mu}.
\end{equation}
For the left edge mode we need $|\lambda| < 1$, so we take $\lambda = \mu/(2t)$
(valid for $|\mu| < 2t$, i.e.\ inside the topological phase).
The wavefunction decays \emph{purely exponentially}:
\begin{equation}
  \psi_j^L \sim \left(\frac{|\mu|}{2t}\right)^{j-1}, \qquad
  \xi = \frac{1}{|\ln|\mu/2t||} \quad \text{(coherence length)}.
\end{equation}
Crucially, for $t = \Delta$ the solution has \textbf{no oscillatory factor}
$e^{ik_0 j}$. This is special to the symmetric point; for $t \neq \Delta$ the
roots can be complex, introducing oscillations.

\subsection{Finite-size hybridization splitting}

A chain of length $L$ has both a left edge mode $\psi^L_j \sim \lambda^{j-1}$
and a right edge mode $\psi^R_j \sim \lambda^{L-j}$.
Their overlap determines the splitting of the zero-energy pair:
\begin{equation}
  \delta E \approx 2t\,\frac{\langle\psi^L | V | \psi^R\rangle}
                             {\|\psi^L\|\,\|\psi^R\|}
  \sim 2t\,\lambda^L = 2t\left(\frac{|\mu|}{2t}\right)^L.
\end{equation}
This is the energy of the near-zero mode.
For $t = \Delta = 1$:
\begin{equation}
  \boxed{E_0(\mu, L) \approx 2\left(\frac{|\mu|}{2}\right)^L.}
  \label{eq:splitting}
\end{equation}
Predictions from \eqref{eq:splitting}:
\begin{itemize}
  \item At the sweet spot $\mu = 0$: $E_0 = 0$ \emph{exactly}, for any $L$.
  \item At the phase boundary $|\mu| = 2t$: $E_0 \to 2$ (the bulk gap), for any $L$.
  \item In the interior: $E_0$ decreases \emph{exponentially} with $L$.
    A larger chain has a better-protected zero mode.
\end{itemize}
So \eqref{eq:splitting} predicts that larger $L$ should give
\emph{smaller} spikes, not larger ones.

\section{Why the spikes increase and widen with $L$: a numerical artifact}

\subsection{The floating-point threshold problem}

The observed behaviour (spikes worsen as $L$ increases) is \emph{not}
a physical effect — it is a numerical artefact arising from how we
define ``positive spectrum''.

In the simulation, positive eigenvalues are selected with the filter
\texttt{evals[evals >= 0]}.  For large $L$, the ideal splitting
$E_0 = 2(|\mu|/2)^L$ falls below machine precision ($\approx 10^{-15}$)
over a wide range of $\mu$:
\begin{equation}
  E_0 < \varepsilon_\text{mach} \iff
  \left(\frac{|\mu|}{2}\right)^L < \frac{\varepsilon_\text{mach}}{2}
  \iff
  |\mu| < 2\,\varepsilon_\text{mach}^{1/L}.
\end{equation}
For $\varepsilon_\text{mach} = 5\times10^{-16}$:
\begin{align}
  L = 30:  &\quad |\mu|_\text{thresh} \approx 0.62, \\
  L = 100: &\quad |\mu|_\text{thresh} \approx 1.41.
\end{align}
For $|\mu| < |\mu|_\text{thresh}$, the near-zero eigenvalue is
indistinguishable from zero at double precision. Due to rounding,
it can be stored as a tiny \emph{negative} number ($\sim -10^{-16}$).
The filter \texttt{evals >= 0} then \emph{excludes} this eigenvalue,
and the \emph{lowest bulk mode} (at finite energy $\sim E_\text{gap}$)
takes its place as the ``lowest positive eigenvalue''. This appears
as a spike.

\subsection{Width and height of the artefact}

\begin{itemize}
  \item \textbf{Width}: the spike region is $|\mu| < 2\varepsilon_\text{mach}^{1/L}$,
    which \emph{grows} with $L$ (for $L=100$ it covers most of the
    topological phase, explaining the ``widened'' spikes).
  \item \textbf{Height}: the spike jumps to the lowest bulk energy,
    which is $\sim E_\text{gap}(\mu) / (L\text{-dependent level spacing})$.
    With larger $L$ the bulk level spacing is smaller, so more bulk
    levels fall inside the window and the spikes look denser.
\end{itemize}

\section{The fix: use quasiparticle energies $|E|$}

The physical quasiparticle spectrum consists of the \emph{absolute values}
of the BdG eigenvalues. The near-zero mode has $|E| \approx 0$ regardless
of the sign of the floating-point result:
\begin{equation}
  \text{correct}: \quad
  \{E_n\}_{n=1}^{L} = \text{sort}\bigl(|e_i|\bigr)_{i=1}^{2L}.
\end{equation}
In code, replace
\begin{verbatim}
  return evals[evals >= 0]          # WRONG: sign-sensitive
\end{verbatim}
with
\begin{verbatim}
  return np.sort(np.abs(evals))[:L]  # correct: quasiparticle energies
\end{verbatim}
This is the standard way to extract the quasiparticle spectrum from a
BdG diagonalisation and is immune to the sign ambiguity at near-zero
eigenvalues.

\section{Summary}

\begin{center}
\renewcommand{\arraystretch}{1.4}
\begin{tabular}{lll}
\hline
Observation & Physical or artefact? & Cause \\
\hline
Spikes near $|\mu|\to 2t$ & \textbf{Physical} & $\xi \to \infty$, hybridisation is strong \\
Spikes in interior for small $L$ & Mostly artefact & Float precision, level crossings \\
Spikes widen/worsen for large $L$ & \textbf{Artefact} & Threshold $|\mu|_\text{thresh}\propto\varepsilon^{1/L}$ grows \\
\hline
\end{tabular}
\end{center}

\end{document}
